\section{Negative heat requires minors of increasing order}
\label{sec:analytic-pf-order-escape}

\begin{lemma}[Uniform angular confinement for bounded negative heat]
\label{lem:analytic-heat-angular}
Fix $M>0$ and $0<\varepsilon<1$. Uniformly for $0\le a\le M$ and
$z$ with $-\operatorname{Im}z\ge\varepsilon|z|$, as $|z|\to\infty$,
\begin{equation}\label{eq:analytic-angular-asymptotic}
 H_{-a}(z)=H_0(z)
   \exp\left(a\left(\frac{H_0'}{H_0}(z)\right)^2\right)
       \left(1+O_{M,\varepsilon}
              \left(\frac{\log^2|z|}{|z|}\right)\right).
\end{equation}
In particular, for sufficiently large $|z|$ this sector contains
no zero of any $H_{-a}$ with $a\in[0,M]$. By conjugation the same
holds in the reflected upper sector. This lemma is unconditional.
\end{lemma}

\begin{proof}
Set $r=|z|$. In the half-plane $\operatorname{Im}w<-1$,
$s=1/2+iw/2$ satisfies $\operatorname{Re}s>1$ and
\[
 H_0(w)=\frac18\xi(s)
       =\frac1{16}s(s-1)\pi^{-s/2}\Gamma(s/2)\zeta(s)
\]
is nonzero. Every factor is nonzero there; in particular there are
no gamma poles and no zeta zeros in this half-plane. It is simply
connected, so choose a holomorphic logarithm $G(w)=\log H_0(w)$.
Write $L=G'(z)$. Stirling's expansion and the absolutely convergent
Euler series in $\operatorname{Re}s>1$ give
\begin{equation}\label{eq:analytic-angular-log-bounds}
 |L|=O_\varepsilon(\log r),\qquad
 G''(w)=O_\varepsilon(r^{-1})\quad
              (|w-z|\le 2r^{1/3}).
\end{equation}
For clarity, the gamma contribution to $G'$ is
$(i/4)\psi(s/2)$ and that to $G''$ is
$-(1/16)\psi'(s/2)$; the polynomial terms are $O(1/r)$ and
$O(1/r^2)$ respectively, and the logarithmic zeta derivatives
are $O(2^{-c_\varepsilon r})$. These estimates are uniform on
the indicated disc for large $r$.

We also need a global bound along the contours used below.
For $p=2aL$ and every real $u$, all points on the straight segment
from $z$ to $z+p+u$ have imaginary part at most
$-\varepsilon r+O_M(\log r)<-\varepsilon r/2$.
Their associated $s$ therefore remain in
$\operatorname{Re}s\ge c_\varepsilon r$. Uniform Stirling estimates
in the right half-plane and the Euler series give
\[
 |G'(w)|\le C_{M,\varepsilon}\log(r+|u|)
\]
along that segment, after enlarging the constant and taking $r$
large. Consequently the holomorphic-logarithm remainder
\[
 R(u)=G(z+p+u)-G(z)-L(p+u)
\]
satisfies
\begin{align}
 |R(u)|&\le
 \frac{C_{M,\varepsilon}}r(|u|+\log r)^2,
                  &&|u|\le r^{1/3},\label{eq:analytic-local-R}\\
 |R(u)|&\le C_{M,\varepsilon}
       (|u|+\log r)\log(r+|u|),
                  &&u\in\mathbb R.\label{eq:analytic-global-R}
\end{align}
The first bound is Taylor's formula with
\eqref{eq:analytic-angular-log-bounds}; the second is integration
of the global derivative bound and $|L|=O(\log r)$.

For $a>0$, the defining Fourier integral gives the exact Gaussian
convolution
\begin{equation}\label{eq:analytic-negative-heat-convolution}
 H_{-a}(z)=\frac1{\sqrt{4\pi a}}
       \int_{\mathbb R}e^{-u^2/(4a)}H_0(z+u)\,du.
\end{equation}
Indeed, integrating $\cos((z+u)v)$ against the Gaussian inserts
$e^{-av^2}$; Fubini is justified by
$|\cos((z+u)v)|\le e^{|\operatorname{Im}z|v}$ for real $u$
and the rapid decay of the theta kernel.

Shift the contour in \eqref{eq:analytic-negative-heat-convolution}
from $\mathbb R$ to $\mathbb R+p$. The integrand is entire in
$u$. Throughout the finite vertical displacement the point $z+u$
still lies in the zero-free half-plane just specified, so the same
branch of $G$ is valid; no gamma pole or zero is crossed.
The Gaussian dominates the growth of $H_0(z+u)$ on the vertical
ends of a large rectangle. More explicitly, for fixed $z,p,a$,
Stirling's bound is at most
$\exp(C_{z,p}|u|\log(2+|u|))$, whereas the Gaussian contributes
$\exp(-(\operatorname{Re}u)^2/(4a)+O_{p,a}(1))$.
Thus those ends tend to zero. With the contour parametrized by
$u+p$, completing the square now gives the exact identity
\[
 \frac{H_{-a}(z)}{H_0(z)e^{aL^2}}
 =\frac1{\sqrt{4\pi a}}
     \int_{\mathbb R}e^{-u^2/(4a)}e^{R(u)}\,du.
\]

The local error is uniform down to $a=0$. On $|u|\le r^{1/3}$,
\eqref{eq:analytic-local-R} tends uniformly to zero and
$|e^{R(u)}-1|\le2|R(u)|$ for large $r$. Its Gaussian mean is
$O_{M,\varepsilon}((a+\log^2r)/r)$.
On $|u|>r^{1/3}$, the right side of
\eqref{eq:analytic-global-R} is at most $u^2/(8M)$ for sufficiently
large $r$, uniformly throughout this tail. Since $a\le M$,
\[
 e^{-u^2/(4a)}|e^{R(u)}|
       \le e^{-u^2/(8a)}.
\]
After normalization, the tail integral is
$O_M(e^{-c_Mr^{2/3}})$, uniformly for $0<a\le M$; the same
bound holds for the omitted Gaussian mass. This proves
\eqref{eq:analytic-angular-asymptotic}. At $a=0$ the assertion is
the identity $H_0=H_0$. Finally, both displayed leading factors
are nonzero, so the relative error proves zero-freeness for
sufficiently large $r$.
\end{proof}

\begin{lemma}[A quadratic factor and the minor-width dictionary]
\label{lem:analytic-quadratic-pfr}
Let $b>0$, $0\le\theta<\pi/(r+1)$, and
\[
 P(w)=(1+be^{i\theta}w)(1+be^{-i\theta}w).
\]
Its coefficient Toeplitz matrix has all minors of order at most
$r$ nonnegative. Products, and locally uniform limits normalized
to have constant coefficient one, retain this property.
\end{lemma}

\begin{proof}
Write $a=be^{i\theta}$, $\bar a=be^{-i\theta}$, so the coefficients
are the elementary symmetric functions in these two variables.
For a $k$ by $k$ minor with increasing row indices $i_1,\ldots,i_k$
and column indices $j_1,\ldots,j_k$, the minor is zero if some
$i_\ell>j_\ell$, by triangularity. Otherwise set
\[
 \lambda'_p=j_{k+1-p}+p-i_1-k,\qquad
 \mu'_p=i_{k+1-p}+p-i_1-k,\qquad 1\le p\le k.
\]
These are partitions with $\mu'\subseteq\lambda'$. Transposing
and reversing both index lists identifies the determinant with
the dual Jacobi--Trudi determinant $s_{\lambda/\mu}(a,\bar a)$.
In particular $\lambda_1\le k$. The bound concerns the width of
the skew diagram, not its possibly arbitrarily large degree.
The dual Jacobi--Trudi/tableau identity follows by the usual
sign-reversing cancellation of intersecting paths in the
determinant expansion; the remaining paths are exactly
semistandard tableaux, with their variable weights.

Here the tableau formula can be evaluated directly, without any
Littlewood--Richardson positivity assumption. If a column of the
skew diagram has more than two boxes, the value is zero. Otherwise
every two-box column must contain $1$ above $2$, contributing
$a\bar a=b^2$. In each row the boxes paired with a box below form
an initial segment, and those paired with a box above form a final
segment, by the partition inequalities. The unpaired boxes form
one intervening interval of length $d\le\lambda_1\le k$. Their
weakly increasing entries can be chosen independently, with
generating function
\[
 h_d(a,\bar a)
   =b^d\frac{\sin((d+1)\theta)}{\sin\theta}>0.
\]
At $\theta=0$ this is $(d+1)b^d$. The full skew Schur value is
therefore a power of $b^2$ times a product of these nonnegative
quantities. This proves the assertion when $k\le r$.

Coefficient convolution is multiplication of the upper Toeplitz
matrices. For each fixed minor the Cauchy--Binet sum is finite,
since all intermediate indices lie between its smallest row and
largest column index. Hence products preserve the inequalities.
Local uniform convergence implies coefficient convergence and
therefore convergence of each fixed finite determinant.
\end{proof}

\begin{theorem}[NN-A020: the detecting minor order escapes]
Under A, for every integer $r\ge1$ there is $\delta_r>0$ such that
the coefficients of
\[
 F_t(w)=H_t(i\sqrt w)/H_t(0)
\]
belong to $PF_r$ for $-\delta_r\le t\le0$. In fact every allowed
minor of order at most $r$, at arbitrary indices, is strictly
positive on this interval. Here \emph{allowed} means
$i_\ell\le j_\ell$ for its increasing index lists.

Let $m_*(t)$ be the minimum order of a negative finite Toeplitz
minor. Under A and Rodgers--Tao, $m_*(t)$ is finite for each $t<0$,
but
\[
                    m_*(t)\longrightarrow\infty
                         \qquad(t\uparrow0).
\]
No simplicity hypothesis on $H_0$ is needed. This concerns minor
order uniformly over all indices; it strengthens the qualitative
finite-window result of NN-A014. The later NN-A022 supplies the
rate $m_*(-a)\gg a^{-1/2}$ by a different uniform product estimate.
\end{theorem}

\begin{proof}
For a zero $z$ in the right half-plane write its angle as
$|\arg z|\in[0,\pi/2)$. There are no zeros on the imaginary axis,
because $H_t(iy)=\int_0^\infty e^{tu^2}\Phi(u)\cosh(yu)\,du>0$.
We first show
\begin{equation}\label{eq:analytic-zero-angles-shrink}
 \sup_{\substack{H_t(z)=0\\\operatorname{Re}z>0}}|\arg z|
                         \longrightarrow0\qquad(t\uparrow0).
\end{equation}
Fix an angle $\eta>0$. Lemma~\ref{lem:analytic-heat-angular}, with
$M=1$ and $\varepsilon=\sin\eta$, excludes all zeros with
$|z|>R_\eta$ and distance in angle at least $\eta$ from the real
axis, uniformly for $-1\le t\le0$. On the remaining compact disc,
$H_t\to H_0$ uniformly. Under A, its finitely many zeros are real
and nonzero. Choose small neighbourhoods of them lying in the
$\eta$ sectors about the positive or negative real axes, and a
zero-free compact complement. Uniform convergence excludes zeros
on that complement for sufficiently small $|t|$. This argument
allows arbitrary multiplicity in the neighbourhoods: it requires
their locations to be real, not their splitting to remain real.
It also excludes a neighbourhood of zero, since $H_0(0)>0$.
This proves \eqref{eq:analytic-zero-angles-shrink}.

Choose $|t|$ small enough that every right-half-plane zero has
$|\arg z|<\pi/(2(r+1))$. The genus-zero product for $F_t$, whose
order is less than one, is absolutely locally convergent and has
constant coefficient one. Its zeros are $w=-z^2$. Real zeros of
$H_t$ give factors $1+bw$ with $b>0$; each nonreal quadruple gives
a quadratic factor as in Lemma~\ref{lem:analytic-quadratic-pfr},
with parameter angle $|\theta|=2|\arg z|<\pi/(r+1)$.
The real linear factors are totally nonnegative. The lemma and
product convergence prove $PF_r$ at all indices.

There are infinitely many factors. Otherwise $F_t$, and hence
$H_t$, would be a polynomial by its order-less-than-one product,
contrary to Fourier decay $H_t(x)\to0$ along the real axis and
$H_t(0)>0$. Every factor has constant term one and strictly
positive linear coefficient. For any allowed minor of order
$k\le r$, move its row configuration to its column configuration
one unit at a time, rightmost particle first, using a different
factor for each move. Each selected single-factor minor is its
positive linear coefficient, while a stationary configuration
has minor one. The particles never collide. This is a positive
term in the finite Cauchy--Binet expansion once sufficiently many
factors are used. Additional factors preserve this lower bound
by their stationary term, and the limiting determinant is
strictly positive.

At $t=0$ the same assertion is NN-A008. Shrinking the chosen
neighbourhood if necessary gives a closed interval
$[-\delta_r,0]$. For each $t<0$, NN-A007, using A and
Rodgers--Tao, says that $F_t$ fails $PF_\infty$, so some finite
minor is negative. But for any fixed $r$ none of orders at most
$r$ is negative near zero. This is precisely the claimed
divergence of $m_*(t)$.
\end{proof}

\begin{corollary}[Effective finite certificates for each fixed order]
Assume A. There is a terminating algorithm which, given $r\ge1$,
returns a positive rational $\delta_r$ and finite analytic bounds
certifying the conclusion of NN-A020. The algorithm does not
need to decide simplicity or the multiplicities of any zero.
No polynomial-time bound is asserted.
\end{corollary}

\begin{proof}
Choose a positive rational $q$, for example $q=1/(4(r+1))$, with
$\arctan q<\pi/(2(r+1))$, and set
$\varepsilon=q/\sqrt{1+q^2}$. The proof of
Lemma~\ref{lem:analytic-heat-angular} is effective: the gamma
estimates have computable remainder bounds and the zeta bounds
come from an absolutely convergent Dirichlet series. Thus its
constants and a rational threshold $R$ for zero-freeness in
$|\operatorname{Im}z|\ge\varepsilon|z|$ can be computed for
$0\le a\le1$. There is no compactness-only constant in that lemma.
For example, the only estimate uniform over an unbounded real
tail is \eqref{eq:analytic-global-R}; after division by $u^2$ its
majorant is
\[
 C\frac{\log(R+|u|)}{|u|}
          \left(1+\frac{\log R}{|u|}\right),
\]
decreasing for $|u|\ge R^{1/3}$ when $R>e$.
The resulting finite elementary inequalities, and the local
$C\log^2R/R$ error, have computable thresholds. Increasing $R$
to a rational value preserves a threshold established by these
eventually decreasing majorants. Explicit digamma and trigamma
bounds follow either from the remainder form of Stirling or
from Euler summation of their absolutely convergent series.

Consider the compact semialgebraic set
\[
 K_{R,q}=\{z:|z|\le R,\ |\operatorname{Im}z|
                                \ge q|\operatorname{Re}z|\}.
\]
Under A, $H_0$ has no zero on this set; the only real point of the
set is zero, where $H_0(0)>0$. Its positive minimum can be
certified without listing its zeros. Cover the set by rational
grid boxes, evaluate $H_0$ at their centers by complex-ball
arithmetic, and subtract a box-radius times a certified derivative
bound on $|z|\le R+1$. Refine the grid and the evaluations until
all these lower bounds are positive. Continuity and the positive
compact minimum guarantee termination. Intersection of a rational
box with $K_{R,q}$ is a real-algebraic decision, or one may use
an effective finite overcover converging to this compact set.
The Fourier integral provides a computable derivative bound.
Denote the resulting rational lower bound by $m_R>0$.

Compute a rational upper bound $B_R$ for
\[
 \int_0^\infty u^2e^{Ru}\Phi(u)\,du.
\]
Such a bound is effective from the explicit theta sum and its
superexponential tail. Since $0\le1-e^{-au^2}\le au^2$,
\[
 \sup_{|z|\le R}|H_{-a}(z)-H_0(z)|\le aB_R.
\]
It follows that
$\delta_r=\min(1,m_R/(2B_R))$ has the required property:
there are no zeros on $K_{R,q}$ for $0\le a\le\delta_r$, and
the angular lemma excludes them on the unbounded part of the
same cones. Every zero therefore lies in
$|\operatorname{Im}z|<q|\operatorname{Re}z|$.
The proved factor criterion and positive paths certify all
orders at most $r$, at arbitrary indices, simultaneously.
\end{proof}

\paragraph{Quantitative interface.}
If
$\alpha(t)=\sup_{H_t(z)=0,\,\operatorname{Re}z>0}|\arg z|$
is bounded explicitly, the same proof certifies $PF_r$ whenever
$2(r+1)\alpha(t)<\pi$. Thus a quantitative angular-confinement
bound would give a quantitative minimum-order lower bound.
The fixed-sector estimate proved here does not control the rate
at which the finitely many bounded zeros approach the real line,
nor its interaction with $R_\eta$ as $\eta\downarrow0$.
It gives neither an exponential collision-height law nor a
lower bound for all algorithms.
The corollary gives an effective fixed-order certificate. The later
NN-A022 strengthens it to the explicit interval
$-2^{-36}/(r+1)^2\le t\le0$ and the corresponding
inverse-square-root order lower bound.

\paragraph{Classical locators and exact geometric scope.}
The sector implication is classical Schoenberg theory. A precise
published statement is Adm--Garloff--Tyaglov, \emph{Total
nonnegativity of finite Hurwitz matrices and root location of
polynomials}, JMAA~467 (2018), 148--170, Theorem~5.8,
equation~(5.10), printed p.~161, and the quadratic case
(5.11)--(5.12), p.~162; DOI~10.1016/j.jmaa.2018.06.065.
The original reference given there is Schoenberg,
\emph{A note on multiply positive sequences and the Descartes
rule of signs}, Rend.\ Circ.\ Mat.\ Palermo~4 (1955), 123--131.
The proof above supplies the needed two-variable assertion
directly; no unverified converse for products is used.
Stirling estimates may be read in DLMF~5.11.1--5.11.2 and
their differentiated form; the actual Gaussian-contour
comparison needed here is proved above.

The Toeplitz minors are the skew Schur/exterior-power tests with
diagram width bounded by their order. The theorem proves that
every fixed such width, even with unbounded polynomial degree or
matrix index, retains positivity for a negative-time interval.
The full infinite Schur cone has nevertheless been left at every
negative time. This is a concrete failure of bounded-order
positivity tests to detect the boundary; no finite-dimensional
pure motive, integral GCT basis, or complexity-class separation
is asserted.
The concrete exterior-power dictionary is SGA~6, Expos\'e~V,
Example~2.2(b) and Lemma~2.2.1, with its representation-ring
example. Those operations supply the symmetric-function
formalism; the actual analytic positivity and width bound
are proved here.
